\documentclass[12pt]{article}
\usepackage[utf8]{inputenc}
\usepackage[T1]{fontenc}
\usepackage{hyperref}
\title{Mariyedi Language Resources for K-12 Curriculum Development}
\author{Compiled by Claude Sonnet 5.5 (claude-sonnet-5-5) for Andr\'as Kornai}
\date{}
\begin{document}
\maketitle

\section{Introduction}
Mariyedi (also Marijedi, Murijadi, Marijadi; ISO 639-3 zmy) is a critically endangered Western Daly language of the northwestern Northern Territory, southwest of Darwin. It is grouped with the Marithiel languages. About twenty speakers were reported in 1981, and today only a few remain. Speakers have shifted to Kriol, and the language has no literacy tradition.

\section{Primary Educational Resources}
None found.

\section{Linguistic and Reference Materials}
\begin{itemize}
\item Endangered Languages Catalogue entry: \url{https://endangeredlanguages.com/lang/2327}
\item Glottolog entry for the Marri Yedi/Marithiel grouping: \url{https://glottolog.org/resource/languoid/id/mari1420}
\item Ian Green's 1989 work on Marrithiyel (PHOIBLE source record): \url{https://phoible.org/sources/green_marrithiyel_1989}
\item WALS reference record, Green 1989: \url{https://wals.info/refdb/record/Green-1989}
\item Grambank entry: \url{https://grambank.clld.org/languages/mari1424}
\end{itemize}

\section{Religious and Community Resources}
None found.

\section{Digital and Audio Resources}
The Daly Languages project is making Ian Green's field notes, recordings and analyses of Daly languages (recorded 1980--1996) accessible to researchers and communities: \url{https://dalylanguages.org/}. Whether it holds Mariyedi material was not verified.

\section{Curriculum Development Context}
There are about six speakers, all elderly. No school programme exists. Material on Mariyedi proper is scarce, since much of the documentation concerns the closely related Marrithiyel. Any curriculum should be led by the speakers and their families, and archival recordings and notes should be consulted with the community's permission.
\end{document}